% SEGMENT OLP-0594-S01
\documentclass[../../../include/open-logic-section]{subfiles}

\begin{document}

\olfileid{sth}{choice}{tarskiscott}
\olsection{டார்ஸ்கி--ஸ்காட் உத்தி}

\olref[cardinals][cardsasords]{defcardinalasordinal} இல்,
கண அளவெண்களை வரிசையெண்களாக வரையறுத்தோம். இதற்காக நன்கு
வரிசைப்படுத்தல் அடிகோளத்தை எடுத்துக்கொண்டோம். அது பொதுவாகப்
பின்பற்றப்படும் வழிமுறை என்பதற்காகவே அவ்வாறு செய்தோம்.

நன்கு வரிசைப்படுத்தலைச் சுற்றியுள்ள தத்துவச் சிக்கல்களைப்
பேசுவதற்கு முன், பொதுவழியிலிருந்து \emph{விலகி}, நன்கு
வரிசைப்படுத்தலை எடுத்துக்கொள்ளாமலேயே கண அளவெண்களின்
கோட்பாட்டை வளர்க்க \emph{முடியும்} என்பதைத் தெளிவுபடுத்த
வேண்டும். \olref[sfr][siz][]{chap} இல் வந்த
$\cardeq{A}{B}$, $\cardle{A}{B}$, $\cardless{A}{B}$
ஆகியவற்றின் வரையறைகளைத் தொடர்ந்து பயன்படுத்தலாம்.
\emph{கண அளவெண்} என்ற கருத்துக்கு மட்டும் புதிய வரையறை தேவை.

முதலில் தோன்றக்கூடிய எளிய எண்ணம், $A$ இன் கண அளவை
பின்வருமாறு வரையறுப்பதாக இருக்கும்:
\[
	\Setabs{x}{\cardeq{A}{x}}.
\]
இதை \olref[cardinals][hp]{sec} இன் இறுதியில் சுருக்கமாகக்
கூறிய ஃபிரேகேயின் $\# x Fx$ வரையறையுடன் ஒப்பிடலாம்.
அங்குச் சுட்டிய காரணங்களாலேயே இந்த வரையறை செயல்படாது.
எந்த ஓருறுப்புக் கணமும் $\{\emptyset\}$ க்கு சமஎண்ணிக்கையுடையது.
ஆனால் படிநிலைமுறையின் ஒவ்வொரு தொடர்ச்சி நிலையிலும் புதிய
ஓருறுப்புக் கணங்கள் தோன்றுகின்றன; முந்தைய படிநிலையின்
ஓருறுப்புக் கணத்தை மட்டும் கருதினாலே போதும். ஆகவே
$\Setabs{x}{\cardeq{A}{x}}$ இருப்பதில்லை; அதற்கு ஒரு
தரநிலையை ஒதுக்க முடியாது.

% SEGMENT OLP-0594-S02
இந்தச் சிக்கலைத் தவிர்க்க டார்ஸ்கியும் ஸ்காட்டும் முன்வைத்த
உத்தியைப் பயன்படுத்துகிறோம்:\footnote{நினைவூட்டல்: வெளிப்படையாக
வேறு நிபந்தனை இல்லையெனில், எல்லா வாய்பாடுகளிலும் அளபுருக்கள்
இருக்கலாம்.}

\begin{defn}[டார்ஸ்கி--ஸ்காட்]
எந்த வாய்பாடு $\phi(x)$ க்கும், $[ x : \phi(x)] $ என்பது
$\phi(x)$ ஐ நிறைவேற்றும் கணங்களுள் சாத்தியமான மீச்சிறிய
தரநிலையுடைய எல்லா $x$ களின் கணம்; $\phi$ ஐ நிறைவேற்றும் பொருள்கள்
இல்லையெனில் அது $\emptyset$.
\end{defn}

இந்த வரையறை செல்லுபடியாகுமா என்று சரிபார்ப்போம். $\ZF$
இல், ஒவ்வொரு $x$ க்கும் $\setrank{x}$ இருக்கிறது என்பதை
\olref[spine][foundation]{zfentailsregularity} உறுதிசெய்கிறது.
இப்போது $\phi$ ஐ நிறைவேற்றும் பொருள்கள் இருந்தால்,
$(\exists x\subseteq V_\alpha)\phi(x)$ நிறைவேறும் மீச்சிறிய
தரநிலையை $\alpha$ என்க; அதாவது
$(\forall \beta \in \alpha)(\forall x \subseteq V_\beta)\lnot \phi(x)$.
அப்போது பிரிப்பின் மூலம் $[x : \phi(x)]$ ஐ
$\Setabs{x \in V_{\alpha+1}}{\phi(x)}$ என்று வரையறுக்கலாம்.

% SEGMENT OLP-0594-S03
டார்ஸ்கி--ஸ்காட் உத்தியை நியாயப்படுத்திய பின்,
அதைக் கொண்டு கண அளவின் ஒரு கருத்தை வரையறுக்கலாம்:

\begin{defn}
$A$ இன் \textsc{ts}-கண அளவு
$\text{tsc}(A) = [x : \cardeq{A}{x}]$.
\end{defn}

\textsc{ts}-கண அளவெண்ணின் வரையறை நன்கு வரிசைப்படுத்தலைப்
பயன்படுத்துவதில்லை. அந்த அடிகோளம் இல்லாமலேயே,
\emph{\textsc{ts}-கண அளவெண்கள்},
\olref[cardinals][cardsasords]{defcardinalasordinal} இல்
வரையறுத்த \emph{கண அளவெண்களைப்} போலவே பல வழிகளில்
செயல்படுகின்றன என்பதை நிறுவலாம். எடுத்துக்காட்டாக,
\olref[cardinals][cardsasords]{lem:CardinalsBehaveRight} மற்றும்
\olref[card-arithmetic][opps]{lem:SizePowerset2Exp}
ஆகியவற்றை \textsc{ts}-கண அளவெண்களைக் கொண்டு மீண்டும்
கூறினால், அவற்றின் நிறுவல்கள் நன்கு வரிசைப்படுத்தலை
எடுத்துக்கொள்ளாமல் $\ZF$ இலேயே செல்லும்.

இதே உத்தியால் வரிசையெண்களின் கோட்பாட்டையும் வளர்க்கலாம்.
$\tuple{A, <}$ ஒரு நன்கு வரிசைப்படுத்தல் எனில்,
$\text{tso}(A, <) = [\tuple{X, R} :
\ordeq{\tuple{A, <}}{\tuple{X, R}}]$ என்க. கண அளவெண்கள்,
வரிசையெண்கள் பற்றிய இந்த அணுகுமுறைக்கு
\citet[chs.~9--12]{Potter2004} ஐப் பார்க்கவும்.

\end{document}
